% Requires booktabs and references.bib. All figures use this profile.
\paragraph{Literature-parameterized physical model.}
We adopt the nominal instruction times reported by Liao et al.\ in
\emph{Benchmarking of quantum protocols}~\cite{liao2022benchmarking}
(Table~3), and the matching H/X/Z, memory and gate-noise parameters
used by Iqbal et al.\ in \emph{Investigating Imperfect Cloning for
Extending Quantum Communication Capabilities}~\cite{iqbal2023cloning}
(Table~3 and Section~4).
\begin{table}[t]
\centering
\caption{Adopted nominal processor profile. Times and noise parameters are
from~\cite{iqbal2023cloning}; instruction times also appear
in~\cite{liao2022benchmarking}. Derived circuit durations are QuCl choices.}
\label{tab:qucl-physical-profile}
\small
\begin{tabular}{@{}lr@{}}
\toprule
Parameter & Value \\
\midrule
CNOT & 20~$\mu$s \\
H / X / Z & 5~ns each \\
Single measurement & 3.7~$\mu$s \\
Memory $T_1$, $T_2$ & 10~h, 1~s \\
Gate depolarization parameter $p$ & 0.01 \\
\midrule
Serial CNOT--H--M--M & 27.405~$\mu$s \\
Two correction slots & 10~ns \\
\bottomrule
\end{tabular}
\end{table}

NetSquid executes the timed instructions and applies native T1/T2 storage
noise. We explicitly implement the quoted gate parameter as independent
per-participating-qubit depolarization,
$\mathcal{D}_p(\rho)=(1-p)\rho+pI/2$, before each actual CNOT, H, X or Z,
after storage noise for the instruction's duration. The channel extends
locally to entangled qubits. Ideal measurements and unused identity
correction slots receive storage noise only. Both correction slots
reserve 5~ns regardless of the measured branch.
This noise placement and the serial measurement schedule are explicit
QuCl modeling choices, not an assertion that the cited simulator uses
the same channel placement. Initialization and readout are ideal;
quantum transit is lossless and noiseless. The profile is a published
nominal abstraction, not a calibrated physical NV device or a complete
reproduction of the cited experiments.
